% !TEX root = ../../thesis.tex
\section{Convergence of the Min--Max Sequence}\label{ch:convergence-of-min-max}

This section completes the proof of the main theorem. Starting from the good sweepouts constructed in the preceding sections, we perform a pull--tight deformation on an exhaustion of~$M$, extract a varifold limit that is almost minimizing in small annuli, and apply the regularity theory of Li--Zhou~\cite{LZ17} to obtain a smooth, almost properly embedded free boundary minimal hypersurface.

Throughout this section, ``stationary'' means stationary in~$M$ with free boundary (Definition~\ref{def:stationary}).

%%%%%%%%%%%%%%%%%%%%%%%%%%%%%%%%%%%%%%%%%%%%%%%%%%%%%%%%%%%%%%%%%%%%%%%%%%%%%%%
\subsection{Pull--tight on a non-compact manifold}\label{sec:pull-tight}
%%%%%%%%%%%%%%%%%%%%%%%%%%%%%%%%%%%%%%%%%%%%%%%%%%%%%%%%%%%%%%%%%%%%%%%%%%%%%%%

Fix a good set $U\subset M$. By Proposition~\ref{prop:nested-sweepout}, there exist good sweepouts of~$U$ with maximal area arbitrarily close to~$W_g(U)$. Since $M$ is non-compact, we cannot pull tight on all of~$M$ at once. Following Chambers--Liokumovich~\cite[\S7]{CL20}, we pull tight layer by layer on a bounded exhaustion and combine the results via a diagonal argument.

Let $U_0\Subset U_1\Subset\cdots\Subset M$ be bounded relative open sets with $M=\bigcup_i U_i$ and $U\Subset U_0$. Since each $\overline{U_i}$ is compact, the monotonicity formula \cite[Theorem~2.1]{LZ17} and Li--Zhou's estimates \cite[Section~3]{LZ17} apply on each~$U_i$. For each~$i$, choose $W_i$ with $U_i\Subset W_i\Subset U_{i+1}$. Set $C:=2W_g(U)$ and define
\[
\mathcal{V}^C(\Omega):=\bigl\{V\in\mathcal{V}_n(M):\operatorname{spt}\|V\|\subset\overline\Omega,\;\|V\|(M)\le C\bigr\},
\]
\[
\mathcal{V}^C_\infty(\Omega):=\bigl\{V\in\mathcal{V}^C(\Omega): V\text{ is stationary in }\Omega\text{ with free boundary}\bigr\}.
\]
Both sets are compact under the varifold metric~$\mathbf{F}$ when $\Omega$ is bounded (see \cite[Lemma~3.4]{LZ17}). For the full (non-compact) manifold, we define analogously
\[
\mathcal{V}^C(M):=\bigl\{V\in\mathcal{V}_n(M): \|V\|(M)\le C\bigr\},
\]
\[
\mathcal{V}^C_\infty(M):=\bigl\{V\in\mathcal{V}^C(M): V\text{ is stationary in }M\text{ with free boundary}\bigr\}.
\]
Note that $\mathcal{V}^C(U_j)\subset\mathcal{V}^C(M)$ for every~$j$, but $\mathcal{V}^C(M)$ itself is \emph{not} compact under~$\mathbf{F}$ when $M$ is non-compact, since a sequence of varifolds whose supports escape to infinity need not have an $\mathbf{F}$-convergent subsequence.

\begin{lemma}[Pull--tight map {\cite[Lemma~2.3]{CL20}}]\label{lem:extend-pull-tight}
For every~$i$ there exist a continuous map $\Phi_{U_i,W_i}:\mathcal{V}^C(M)\to\mathcal{V}^C(M)$ with the following properties for all $V\in\mathcal{V}^C(M)$:
\begin{enumerate}
    \item \textup{(Mass non-increase)} $\|\Phi_{U_i,W_i}(V)\|(M)\le \|V\|(M)$;
    \item \textup{(Inactivity below threshold)} If $\|V\|(\overline{W}_i)\le 5\,\mathbf{M}(\partial_IU)$, then $\Phi_{U_i,W_i}(V)=V$;
    \item \textup{(Strict decrease for non-stationary varifolds)} If $\|V\|(\overline{W}_i)\ge 9\,\mathbf{M}(\partial_IU)$ and $\mathbf{F}_{U_i}\bigl(V\llcorner \mathrm{Gr}_n(U_i),\,\mathcal{V}^C_\infty(U_i)\bigr)\in\bigl[2^{-(k+1)},2^{-k}\bigr]$ for some $k\in\mathbb{N}$, then:
    \begin{enumerate}[label=\textup{(\alph*)}]
        \item $\|\Phi_{U_i,W_i}(V)\|(M) \le \|V\|(M) - \varepsilon_k(U_i)$;
        \item $\mathbf{F}\bigl(V,\,\Phi_{U_i,W_i}(V)\bigr)\le\tau_k(U_i)$;
    \end{enumerate}
    where $\varepsilon_k(U_i)\to 0$ and $\tau_k(U_i)\to 0$ as $k\to\infty$.
\end{enumerate}
\end{lemma}

\begin{proof}
The construction follows \cite[Lemma~2.3]{CL20} verbatim, with the sole modification that all vector fields are taken in~$\mathfrak{X}_{\mathrm{tan}}(M)$ (tangent to~$\partial M$) rather than~$\mathfrak{X}_c(M)$. This does not affect the argument: the proof uses only that $\mathfrak{X}_{\mathrm{tan}}(M)$ is a closed subspace of~$\mathfrak{X}_c(M)$ on which the covering, partition-of-unity, and flow constructions remain valid, and that non-stationarity with free boundary provides a vector field $X\in\mathfrak{X}_{\mathrm{tan}}(M)$ with $\delta V(X)<0$.
\end{proof}

\begin{remark}[Role of the thresholds]\label{rem:thresholds}
The constants $5\,\mathbf{M}(\partial_IU)$ and $9\,\mathbf{M}(\partial_IU)$ in Lemma~\ref{lem:extend-pull-tight} are not explained in~\cite{CL20} but play an important role. Endpoints of a good sweepout have mass at most $5\,\mathbf{M}(\partial_IU)$ (Definition~\ref{def:good-sweepout}), so property~(2) leaves them unchanged. Moreover, since the flow is generated by vector fields supported in~$W_i$ with $U\Subset U_0\Subset W_0$, it does not transport mass across~$\partial_IU$. Together, these guarantee that the pull-tight preserves the good sweepout structure, so the non-escape property (Proposition~\ref{prop:non-escape}) remains applicable. On the other hand, the min--max slices have mass close to $W_g(U)\ge\tfrac{9}{10}\,W_\partial(U)\gg 9\,\mathbf{M}(\partial_IU)$ (by the good set condition), so they are well above the threshold in property~(3) and are actively pulled tight. The gap between $5$ and $9$ accommodates a continuous cutoff.
\end{remark}

Let $\{\Gamma_t^i\}$ be a sequence of good sweepouts with $\sup_t \mathbf{M}(\Gamma_t^i)\downarrow W_g(U)$. On a compact manifold, one would apply a single pull-tight map to the entire sequence. Here $M$ is non-compact, so instead we pull tight incrementally: the $i$-th sweepout is pulled tight on the first $i$ layers of the exhaustion. Set
\[
\mathbf{P}_i:=\Phi_{U_0,W_0}\circ\Phi_{U_1,W_1}\circ\cdots\circ\Phi_{U_i,W_i}.
\]
Thus $\mathbf{P}_i(\Gamma_t^i)$ has been pulled tight on $U_0,\ldots,U_i$: as $i\to\infty$ the number of layers grows with the exhaustion, so the limit varifold will be stationary on every~$U_j$ and hence on all of~$M$.

\begin{proposition}[Pulled--tight minimizing sequence]\label{prop:pt-minimizing}
After passing to a subsequence, $\{\mathbf{P}_i(\Gamma_t^i)\}$ is a minimizing sequence of good sweepouts with
\[
\lim_{i\to\infty}\ \sup_{t\in[0,1]}\mathbf{M}\bigl(\mathbf{P}_i(\Gamma_t^i)\bigr)=W_g(U),
\]
and any min--max sequence $\mathbf{P}_i(\Gamma_{t_i}^i)$ admits a subsequence converging to a varifold $V\in\mathcal{V}^C_\infty(M)$ that is stationary with free boundary.
\end{proposition}

\begin{proof}
The free boundary condition does not affect the argument: it relies only on the compactness of~$\mathcal{V}^C(U_j)$ and on properties~(1)--(3) of Lemma~\ref{lem:extend-pull-tight}, all of which hold because the vector fields lie in~$\mathfrak{X}_{\mathrm{tan}}(M)$. The proof follows~\cite[\S7]{CL20}.

First, a contradiction argument using the mass decrease and $\mathbf{F}$-displacement bounds in Lemma~\ref{lem:extend-pull-tight}(3) shows that for every~$i$,
\begin{equation}\label{eq:pt-claim}
\lim_{j\to\infty}\mathbf{F}_{U_i}\bigl(\mathbf{P}_j(\Gamma_{t_j}^j)\llcorner\mathrm{Gr}_n(U_i),\;\mathcal{V}^C_\infty(U_i)\bigr)=0:
\end{equation}
if the $\mathbf{F}_{U_i}$-distance stayed bounded away from zero, then either some intermediate layer $\Phi_{U_l,W_l}$ would strictly decrease the mass by the $\varepsilon_k$-bound, or the $\tau_k$-bound and the triangle inequality would force the final varifold close to $\mathcal{V}^C_\infty(U_i)$; in both cases one reaches a contradiction with $\mathbf{M}(\mathbf{P}_j(\Gamma_{t_j}^j))\to W_g(U)$. See~\cite[\S7]{CL20} for the full details.

Next, a diagonal argument produces a subsequence converging in $\mathbf{F}_{U_j}$ for every~$j$. Since $M = \bigcup_j U_j$, the limit varifold $V$ is well-defined on all of~$M$, and by~\eqref{eq:pt-claim} we have $V\llcorner\mathrm{Gr}_n(U_j)\in\mathcal{V}^C_\infty(U_j)$ for every~$j$. It follows that $V$ is stationary with free boundary on~$M$: for any $X\in\mathfrak{X}_{\mathrm{tan}}(M)$, the support of~$X$ is contained in some~$U_j$, and $V$ is stationary on~$U_j$, so $\delta V(X) = 0$.
\end{proof}

%%%%%%%%%%%%%%%%%%%%%%%%%%%%%%%%%%%%%%%%%%%%%%%%%%%%%%%%%%%%%%%%%%%%%%%%%%%%%%%
\subsection{Almost minimizing varifolds with free boundary}\label{sec:am-varifolds}
%%%%%%%%%%%%%%%%%%%%%%%%%%%%%%%%%%%%%%%%%%%%%%%%%%%%%%%%%%%%%%%%%%%%%%%%%%%%%%%

The pulled-tight varifold $V$ from Proposition~\ref{prop:pt-minimizing} is stationary, but stationarity alone does not yield regularity. Following the Almgren--Pitts strategy, one needs the stronger \emph{almost minimizing} property: roughly, $V$ cannot be deformed to decrease its mass by much. In the closed (boundaryless) setting, this notion was introduced by Pitts~\cite{Pit81} and reformulated by De~Lellis--Tasnady~\cite{DLT13}; the regularity of almost minimizing varifolds was established in~\cite{Pit81,DLT13}. Chambers--Liokumovich~\cite{CL20} use these results directly, as their manifolds have no boundary. For manifolds \emph{with} boundary, Li--Zhou~\cite[Sections~4--6]{LZ17} developed free boundary counterparts of both the almost minimizing definitions and the regularity theory, which are the versions we use here.

In this section we collect the relevant definitions and results from~\cite{LZ17}. The definitions below are local and apply to any complete manifold with boundary. Of the two main results that use them, the regularity theorem \cite[Theorem~6.2]{LZ17} has an entirely local proof (replacements, tangent cones, and boundary blow-ups all take place in a neighbourhood of a single point) and therefore extends to non-compact manifolds without modification. The only step requiring compactness of the ambient manifold is the combinatorial selection of almost minimizing varifolds \cite[Theorem~5.3]{LZ17}; in our setting this is handled by applying the argument on each bounded set in the exhaustion and passing to a diagonal subsequence, exactly as in~\cite[\S7]{CL20}. The details are carried out in Section~\ref{sec:main-proof}.

The basic notion captures cycles whose mass cannot be decreased much by local deformations supported in a given open set.

\begin{definition}[Almost minimizing with free boundary {\cite[Definition~4.18]{LZ17}}]\label{def:am-free-bdy}
Let $U\subset M$ be relatively open. For $\varepsilon,\delta>0$, define $\mathscr{A}_n(U;\varepsilon,\delta;\mathcal{F})$ to be the set of all $\tau\in\mathbf{Z}_n(M,\partial M;\mathbb{Z}_2)$ such that whenever $\tau=\tau_0,\tau_1,\ldots,\tau_m\in\mathbf{Z}_n(M,\partial M;\mathbb{Z}_2)$ is a discrete sequence with
\begin{itemize}
    \item $\operatorname{spt}(\tau_i-\tau)\subset U$,
    \item $\mathcal{F}(\tau_{i+1}-\tau_i)\le\delta$,
    \item $\mathbf{M}(\tau_i)\le\mathbf{M}(\tau)+\delta$ for $i=1,\ldots,m$,
\end{itemize}
then $\mathbf{M}(\tau_m)\ge\mathbf{M}(\tau)-\varepsilon$.

A varifold $V\in\mathcal{V}_n(M)$ is \emph{almost minimizing in $U$ with free boundary} if there exist sequences $\varepsilon_i\to 0$, $\delta_i\to 0$, and $\tau_i\in\mathscr{A}_n(U;\varepsilon_i,\delta_i;\mathcal{F})$ with canonical representatives $T_i\in\tau_i$ such that $\mathbf{F}(|T_i|,V)\le\varepsilon_i$.
\end{definition}

For regularity, one needs the almost minimizing property to hold not just in a single open set but in sufficiently small annuli around every point; this ensures that every tangent cone is area-minimizing. At interior points, the annuli are ordinary; near the boundary, one uses half-annuli.

\begin{definition}[Almost minimizing in small annuli {\cite[Definition~5.1]{LZ17}}]\label{def:am-small-annuli}
For $p\in M$ and $0<s<r$, set
\[
\mathcal{A}_{s,r}(p):=
\begin{cases}
B_r(p)\setminus\overline{B}_s(p) & \text{if } p\in\operatorname{int}(M),\\
B_r^+(p)\setminus\overline{B}^+_s(p) & \text{if } p\in\partial M,
\end{cases}
\]
where $B_r^+(p):=B_r(p)\cap M$. A varifold $V\in\mathcal{V}_n(M)$ is \emph{almost minimizing in small annuli with free boundary} if for each $p\in M$ there exists $r_{\mathrm{am}}(p)>0$ such that $V$ is almost minimizing in $\mathcal{A}_{s,r}(p)$ with free boundary for all $0<s<r\le r_{\mathrm{am}}(p)$.
\end{definition}

A first consequence is that the almost minimizing property already implies stationarity, so no separate stationarity assumption is needed in the regularity theorem.

\begin{proposition}[{\cite[Proposition~4.21]{LZ17}}]\label{prop:am-stationary}
If $V\in\mathcal{V}_n(M)$ is almost minimizing in a relatively open set $U\subset M$ with free boundary, then $V$ is stationary in $U$ with free boundary.
\end{proposition}

The main payoff is that almost minimizing in small annuli, combined with stationarity, gives full regularity. This is the deepest result in~\cite{LZ17}, relying on replacements, tangent cone classification, and boundary blow-up analysis.

\begin{theorem}[Regularity {\cite[Theorem~6.2]{LZ17}}]\label{thm:main-regularity}
Let $2\le n\le 6$. Suppose $V\in\mathcal{V}_n(M)$ is stationary in $M$ with free boundary and almost minimizing in small annuli with free boundary. Then there exist $N\in\mathbb{N}$, positive integers $m_1,\ldots,m_N$, and smooth, compact, connected, almost properly embedded free boundary minimal hypersurfaces $(\Sigma_j,\partial\Sigma_j)\subset(M,\partial M)$ such that
\[
V=\sum_{j=1}^N m_j\,|\Sigma_j|.
\]
The restriction $2 \leq n \leq 6$ ensures that the singular set is empty; for $n \geq 7$ it has Hausdorff dimension at most $n - 7$.
\end{theorem}

%%%%%%%%%%%%%%%%%%%%%%%%%%%%%%%%%%%%%%%%%%%%%%%%%%%%%%%%%%%%%%%%%%%%%%%%%%%%%%%
\subsection{Proof of the main theorem}\label{sec:main-proof}
%%%%%%%%%%%%%%%%%%%%%%%%%%%%%%%%%%%%%%%%%%%%%%%%%%%%%%%%%%%%%%%%%%%%%%%%%%%%%%%

We now assemble the ingredients. Let $V\in\mathcal{V}^C_\infty(M)$ be the pulled-tight varifold limit from Proposition~\ref{prop:pt-minimizing}. In the compact setting, Li--Zhou~\cite[Theorem~5.3]{LZ17} obtain the almost minimizing property from the combinatorial selection argument of Almgren--Pitts~\cite{Pit81}. The key observation for the non-compact case is that this argument is \emph{local}: for each $p\in M$, one only deforms sweepouts inside a bounded neighbourhood of~$p$. Since each~$U_j$ in the exhaustion is bounded, the combinatorial argument applies on~$U_j$ without modification, and a diagonal argument over~$j$ yields:

\begin{proposition}[Almost minimizing selection]\label{prop:am-selection}
After passing to a further subsequence, $V$ is almost minimizing in small annuli with free boundary (Definition~\ref{def:am-small-annuli}).
\end{proposition}

Combining Proposition~\ref{prop:am-selection} with the regularity theorem (Theorem~\ref{thm:main-regularity}), we obtain smooth structure. As noted in Section~\ref{sec:am-varifolds}, the proof of Theorem~\ref{thm:main-regularity} is entirely local, so it applies to our non-compact~$M$ without modification:

\begin{proposition}[Regularity of the limit]\label{prop:regular-limit}
There exist $N\in\mathbb{N}$, positive integers $m_1,\ldots,m_N$, and smooth, compact, connected, almost properly embedded free boundary minimal hypersurfaces $(\Sigma_j,\partial\Sigma_j)\subset (M,\partial M)$ such that $V=\sum_{j=1}^N m_j\,|\Sigma_j|$.
\end{proposition}

With these preparations, we restate and prove the main theorem (Theorem~\ref{thm:main}).

\begin{theorem}[Main Theorem]\label{thm:main-complete}
Let $(M^{n+1}, g)$ be a complete Riemannian manifold with non-empty smooth boundary $\partial M$, $2 \leq n \leq 6$.  Let $U \subset M$ be a \emph{good set} (Definition~\ref{def:good-set}), i.e., a bounded open set of finite perimeter satisfying $\mathbf{M}(\partial_I U) \leq \frac{1}{10}\, W_\partial(U)$, with $\overline{U} \cap \partial M \neq \varnothing$.  Then there exists a smooth, almost properly embedded minimal hypersurface $(\Sigma,\partial\Sigma)\subset(M,\partial M)$ with (possibly empty) boundary $\partial\Sigma\subset\partial M$ such that
\begin{enumerate}[label=\textup{(\roman*)}]
    \item $\Sigma$ is minimal in the interior and meets $\partial M$ orthogonally along $\partial\Sigma$;
    \item \textup{(Area bound)} $\mathcal{H}^n(\Sigma)\le W_\partial(U)+\mathbf{M}(\partial_IU)$;
    \item \textup{(Localization)} $\Sigma\cap\overline{U}\neq\varnothing$.
\end{enumerate}
\end{theorem}

\begin{proof}
The proof combines all preceding chapters in three steps.

\medskip
\noindent\textbf{Step 1: Min--max sequence.}
The good set condition gives $W_g(U) > 0$: indeed, $\mathbf{M}(\partial_I U) \leq \frac{1}{10} W_\partial(U)$ implies $W_\partial(U) > 0$, and Proposition~\ref{prop:width-inequalities} gives $W_g(U) \geq W_\partial(U) - \mathbf{M}(\partial_I U) \geq \frac{9}{10} W_\partial(U) > 0$.

By Proposition~\ref{prop:nested-sweepout}, there exist good sweepouts $\{\Gamma_t^i\}$ of $U$ with $\sup_t\mathbf{M}(\Gamma_t^i)\downarrow W_g(U)$. The pull--tight procedure (Section~\ref{sec:pull-tight}) produces a minimizing sequence $\{\mathbf{P}_i(\Gamma_t^i)\}$ whose varifold limit $V$ is stationary in $M$ with free boundary (Proposition~\ref{prop:pt-minimizing}).

\medskip
\noindent\textbf{Step 2: Regularity.}
By Propositions~\ref{prop:am-selection} and~\ref{prop:regular-limit},
\[
V=\sum_{j=1}^N m_j\,|\Sigma_j|,
\]
where each $(\Sigma_j,\partial\Sigma_j)\subset(M,\partial M)$ is a smooth, compact, connected, almost properly embedded minimal hypersurface meeting $\partial M$ orthogonally along~$\partial\Sigma_j$. Set $\Sigma:=\bigcup_{j=1}^N\Sigma_j$.

\medskip
\noindent\textbf{Step 3: Area bound and localization.}
The area bound follows from the width inequalities (Proposition~\ref{prop:width-inequalities}):
\[
\mathcal{H}^n(\Sigma)\le\|V\|(M)=W_g(U)\le W_\partial(U)+\mathbf{M}(\partial_IU).
\]
For the localization, the non-escape estimate (Proposition~\ref{prop:non-escape}) guarantees $\mathcal{H}^n(\Sigma\cap \overline{U})\ge\varepsilon(U)>0$, so in particular $\Sigma\cap\overline{U}\neq\varnothing$.
\end{proof}

As an immediate corollary, we obtain existence under a volume growth condition.

\begin{corollary}\label{cor:sublinear-existence}
Let $(M^{n+1},g)$ be a complete Riemannian manifold with non-empty smooth boundary $\partial M$, $2\le n\le 6$, and sublinear volume growth:
\[
\liminf_{R\to\infty}\frac{\mathcal{H}^{n+1}(B_R(x)\cap M)}{R}=0\qquad\text{for some }x\in M.
\]
Suppose in addition that there exists a good set $U$ with $\overline{U}\cap\partial M\neq\varnothing$ (e.g., when $\partial M$ is non-compact). Then $M$ contains a smooth, almost properly embedded minimal hypersurface $(\Sigma,\partial\Sigma)\subset(M,\partial M)$ satisfying conclusions~\textup{(i)--(iii)} of Theorem~\ref{thm:main-complete}.
\end{corollary}

\begin{proof}
By Lemma~\ref{lem:good-set-existence}, the sublinear growth condition guarantees the existence of a good set $U\subset M$ with $\overline{U}\cap\partial M\neq\varnothing$. The result follows from Theorem~\ref{thm:main-complete}.
\end{proof}
